\chapter{Introducción}
\label{cap:introduccion}

\section{Motivación}

La monitorización de la salud de motores aeronáuticos (\emph{Engine Health Monitoring}, EHM)
sostiene hoy decisiones de mantenimiento que mueven cifras enormes: un \emph{shop visit} de un
turbofan comercial cuesta del orden de millones de dólares, y una retirada no programada
multiplica ese coste por la disrupción operativa. La técnica de referencia desde los años setenta
es el \emph{Gas Path Analysis} (GPA)~\cite{Urban1975}: inferir el estado de salud de los
componentes del camino de gas (fan, compresores, turbinas) a partir de las medidas disponibles a
bordo (velocidades de rotor, temperatura de gases de escape, flujo de combustible, presiones).
Sobre el GPA se construyó toda una generación de sistemas industriales ---filtros de Kalman,
líneas base termodinámicas, reglas expertas de aislamiento~\cite{Volponi2014}--- que aquí
llamamos \emph{EHM tradicional}.

En la última década, el aprendizaje automático ha producido resultados llamativos en pronóstico
de vida útil remanente (RUL) sobre bancos de datos sintéticos como
C-MAPSS~\cite{Saxena2008} y N-CMAPSS~\cite{AriasChao2021}. Sin embargo, la literatura adolece de
un problema estructural: \textbf{casi nunca se compara la IA contra un enfoque tradicional
implementado con rigor}, sobre los mismos datos, con las mismas métricas y con el mismo esfuerzo
de ajuste. El resultado son comparaciones contra ``hombres de paja'' que impiden saber qué aporta
realmente la IA, cuándo y a qué coste.

\section{Objetivo}

Este trabajo construye un banco de pruebas reproducible en el que ambas familias de EHM
---tradicional y basada en IA--- observan \emph{exactamente los mismos datos} de una flota
sintética del turbofan CFM56-7B26 y se evalúan con \emph{exactamente las mismas métricas}. La
clave metodológica es que los datos sintéticos incorporan \textbf{verdad-terreno completa}: el
estado real de salud de cada componente (desviaciones de eficiencia y capacidad de flujo) en cada
vuelo. Esto, imposible con datos reales de aerolínea, permite medir no solo si un método detecta
o pronostica bien, sino si \emph{atribuye el daño al componente correcto} ---el talón de Aquiles
del GPA clásico (el fenómeno de \emph{smearing}) y la caja negra de la IA.

Al no existir ni un modelo de prestaciones del motor ni datos de degradación disponibles, el
proyecto construye ambos: un gemelo termodinámico del CFM56-7B26 en
pyCycle~\cite{Hendricks2019} calibrado contra datos públicos certificados
(\cref{cap:modelo}), y un generador de flota degradada con firmas de fallo físicas.

\section{Hipótesis}
\label{sec:hipotesis}

Las hipótesis se formalizan con umbrales cuantitativos \emph{antes} de ver los resultados del
conjunto de test (protocolo pre-registrado, \cref{sec:protocolo}). Todas se contrastan con test
de Wilcoxon pareado por motor, corrección de Holm ($\alpha = 0{,}05$) e intervalos bootstrap BCa.

\begin{table}[htbp]
  \centering\small
  \caption{Hipótesis del proyecto y criterios de confirmación.}
  \label{tab:hipotesis}
  \begin{tabular}{p{0.6cm}p{5.4cm}p{7.6cm}}
    \toprule
    ID & Hipótesis & Criterio de confirmación \\
    \midrule
    H1 & La IA detecta antes y con menos falsas alarmas que el \emph{trend monitoring} &
      \emph{lead time} mediano $\geq +20$\,\% y FPR $\leq 0{,}5\times$ frente al trending, a
      \emph{recall} fijo $0{,}9$; $p<0{,}05$ \\
    H2 & La IA aísla mejor fallos con firmas solapadas &
      $+10$ puntos de exactitud de aislamiento en el subconjunto ``confusable'' (firmas a
      $<15^{\circ}$ en el espacio de la ICM) frente a GPA-WLS; $p<0{,}05$ \\
    H3 & La IA pronostica mejor el RUL &
      mejora simultánea en RMSE y en el \emph{score} asimétrico NASA PHM08; $p<0{,}05$ \\
    H4 & El híbrido \emph{physics-informed} domina con pocos datos &
      RMSE de RUL $\leq$ IA pura con el 100\,\% del entrenamiento y estrictamente mejor con el
      10\,\% y el 25\,\%; $p<0{,}05$ \\
    H5 & Los intervalos \emph{conformal} calibran mejor que la covarianza del Kalman &
      menor $|$cobertura empírica $-$ 0{,}9$|$ con anchura media no peor que $+20$\,\% \\
    \bottomrule
  \end{tabular}
\end{table}

La refutación de cualquier hipótesis es también un resultado publicable: el pre-registro la hace
creíble. Queda prohibido reformular hipótesis tras la etiqueta \codigo{prereg-v1} del
repositorio.

\section{Contribuciones}
\label{sec:contribuciones}

\begin{enumerate}
  \item[\textbf{C1.}] \textbf{SynCFM56}: primer conjunto de datos sintético público de GPA sobre
    un motor comercial concreto (CFM56-7B26), con verdad-terreno de parámetros de salud por
    vuelo y formato de \emph{snapshots} tipo ACARS (informe de despegue + informe de crucero por
    vuelo), más realista que las series continuas de C-MAPSS.
  \item[\textbf{C2.}] Comparación cabeza a cabeza con \textbf{protocolo pre-registrado} e igual
    presupuesto de ajuste (mismo número de \emph{trials} de Optuna) para ambas familias.
  \item[\textbf{C3.}] Híbrido \emph{physics-informed} con tres mecanismos ablacionados por
    separado: residuos contra el gemelo digital, pérdidas con restricciones físicas y
    \emph{stacking} GPA$\rightarrow$ML.
  \item[\textbf{C4.}] Cuantificación de incertidumbre comparada: intervalos de RUL por
    \emph{conformal prediction} frente a la covarianza del filtro de Kalman.
  \item[\textbf{C5.}] \textbf{Physics-Consistency Score (PCS)}: métrica nueva de explicabilidad
    que proyecta las atribuciones SHAP al espacio físico de parámetros de salud vía la matriz de
    coeficientes de influencia.
  \item[\textbf{C6.}] Mapa de determinabilidad: ablación sensores $\times$ ruido $\times$ tamaño
    de datos que muestra \emph{cuándo} gana cada enfoque.
  \item[\textbf{C7.}] Validación \emph{sim-to-real}: la metodología completa se re-ejecuta sobre
    N-CMAPSS para comprobar que el \emph{ranking} de métodos no es un artefacto del simulador
    propio.
\end{enumerate}

\section{Estructura del documento}

El \cref{cap:fundamentos} fija los fundamentos: el motor, la formulación del GPA y la física de
la degradación. El \cref{cap:metodologia} describe la arquitectura del banco de pruebas y las
salvaguardas metodológicas. El \cref{cap:modelo} documenta el gemelo termodinámico del
CFM56-7B26 y su calibración contra datos certificados ---el trabajo completado hasta la fecha---
con las evidencias cuantitativas del proceso. El \cref{cap:trabajo-en-curso} resume las fases
restantes y sus criterios de aceptación.
